\documentclass[10pt,a4paper]{amsart}
\usepackage[margin=19mm]{geometry}
\usepackage{amsmath,amssymb,mathtools}
\usepackage[scr=rsfs,cal=euler]{mathalfa}
\usepackage{xfrac}
\usepackage[unicode,hidelinks,bookmarks=false]{hyperref}
\newcommand{\ie}{i.e.\ }
\newcommand{\cf}{cf.\ }
\newcommand{\resp}{respectively\ }
\newcommand{\Subsection}{\subsection}
\providecommand{\nobreakdash}{\nobreak\hbox{-}}
\setlength{\emergencystretch}{2em}
\multlinegap0pt
\usepackage{xcolor}
\usepackage[shortlabels]{enumitem}
\usepackage{mathtools}
\newtheorem{lemma}{Lemma}
\title{\Large\bf The Stochastic Team Orienteering Problem}
\author{Alberto Guastalla, Jean-François Côté, Roberto Aringhieri}
\date{Source: arXiv:2608.29119v2 (CC BY 4.0)}
\begin{document}
\maketitle

\noindent\emph{Source and licence.} \Large\bf The Stochastic Team Orienteering Problem. Version arXiv:2608.29119v2 (CC BY 4.0), distributed by arXiv under the Creative Commons Attribution 4.0 International licence, http://creativecommons.org/licenses/by/4.0/, as recorded in the arXiv licence metadata for this exact version and shown on the abstract page https://arxiv.org/abs/2608.29119v2. Full licence notice: \texttt{licenses/arxiv-2608.29119-CC-BY-4.0.txt}.\par\smallskip

\noindent\textit{Editorial scope.} Counted unit: the article's lemma lem:lemma6 (source0.tex lines 833--836). The article's complete original proof of it is delivered with it (source0.tex lines 838--872, one \texttt{proof} environment of 35 lines). No mathematics is written, repaired or re-derived: the statement and the proof are the article's own lines, in the article's own notation, and no retained line is substituted or re-flowed.  No further source-local context is needed: every cross-reference of every retained line resolves inside the two delivered spans.  Nothing was omitted from any retained span, so the package carries no omission record.  No retained line contains a citation, so the wrapper carries no bibliography and no bibliography entry.  No retained line contains a figure environment, an image-inclusion command or drawing code, so no image is delivered and the package declares no asset dependency.  Editorial formatting only, in addition: an AMS \texttt{amsart} preamble that restates the article's own declarations and macros used by the retained lines, namely the environments \texttt{lemma} on separate counters, together with the standard packages those lines need.  The retained environments print in this document as Lemma 1 in order of appearance, whereas the article prints the same items with the numbering of its own full document.  The complete native source of the article as distributed in the arXiv e-print for this exact version is bundled unchanged as \texttt{sources/208829/source0.tex}.
\begin{lemma}
\label{lem:lemma6}
$\mu_{\text{UB}}$ is a valid upper bound on the route expectation.
\end{lemma}

\begin{proof}
By definition, $\mu_{\text{UB}}$ represents the (${1-\alpha}$)-quantile of the Normal distribution $\bar{\xi} \sim \mathcal{N}(T_{\max}, \sigma^2_{\text{LB}})$:
%
\[
\mathbb{P}(\bar{\xi} \leq \mu_{\text{UB}}) = 1 - \alpha.
\]
%
Since the normal distribution is symmetric with respect to its mean, we have:
%
\[
\mathbb{P}(\bar{\xi} \leq 2 \, T_{\max} - \mu_{\text{UB}}) = \alpha.
\]
%
Subtracting the quantity $T_{\max} - \mu_{\text{UB}}$ from both sides inside the probability term, it follows that:
\begin{equation}
\label{eq:quantile}
\mathbb{P}(\bar{\xi} - (T_{\max} - \mu_{\text{UB}}) \leq T_{\max}) = \alpha,
\end{equation}
where the mean of the Normal random variable $\bar{\xi} - (T_{\max} - \mu_{\text{UB}})$ is equal to $\mu_{\text{UB}}$.
Therefore, for any route $r$ with a route uncertain travel time $\xi_r \sim \mathcal{N}(\mu_r, \sigma^2_{\text{LB}})$ such that $\mu_r \leq \mu_{\text{UB}}$, it holds that:
%
\[
\delta_r \geq \alpha,
\]
%
which means that $r$ is feasible.
%
Thus, the probability $\delta_r$ achieves the minimum value when $\mu_r = \mu_{\text{UB}}$, where $\delta_r = \alpha$ as in~\eqref{eq:quantile}. Consequently, there is no route $r$ that satisfies:
%
\begin{align*}
\mu_r > \mu_{\text{UB}} \, \text{ and } \delta_r \geq \alpha,
\end{align*}
%
considering a route variance equal to $\sigma^2_{\text{LB}}$. This implies that $\mu_{\text{UB}}$ represents a valid upper bound on the route expectation.
\end{proof}

\end{document}
